\documentclass{article}
\usepackage[utf8]{inputenc}
\usepackage[T1]{fontenc}
\usepackage[margin=2cm]{geometry}
\usepackage{hyperref}
\sloppy
\begin{document}
\section*{Historical source}
Encyclopaedia Britannica scan, source row 473. Parse 3.5 extracted the text. Super 3.5 VL corrected it from an enlarged crop of each block. It can still contain recognition errors.
Page headers, footers, and page numbers are removed.

\subsection*{CONCHOLOGY (E0022)}
CONCHOLOGY.

Chap. 1

Some species of patella, which were brought from Madeira, were also subjected to chemical examination, by the same philosopher. When exposed to a red heat in a crucible, they emitted a perceptible smell of horn or feathers; and by farther examination, by solution, the proportion of carbonic matter depleted appeared to be greater, and the proportion of carbonate of lime less, than what was indicated by the result of the experiments on porcellaneous shells. When unburnt shells belonging to the same species, were immersed in nitric acid very much diluted, the epidermis separated, and the whole of the carbonate of lime was dissolved. A gelatinous substance, nearly in a liquid state, remained, but it did not retain the figure of the shell, and exhibited no appearance of a fibrous structure. These shells, therefore, contain a larger portion of gelatinous matter than the porcellaneous shells, but the other component part consists entirely of carbonate of lime.

69
Mothers of pearl.

Shells composed of mother-of-pearl.—Shells of this description were subjected to similar experiments with the former. When the common oyster was exposed to a red heat, the effects were the same as those which were produced by the same process on the species of patella from Madeira. The solution of the unburnt shell was also similar, excepting only that the gelatinous part was of a greater consistency. When the river mussel was burnt in a crucible, it emitted much smoke, with a strong smell of burnt horn or cartilage; the shell became of a dark gray colour, and exfoliated. By solution in the acids, the proportion of carbonic matter separated was greater, and that of carbonate of lime obtained was less, than from the other shells on which experiments were made.

When an unburnt shell of this description was immersed in diluted nitric acid, a rapid solution and effervescence took place; and at the end of two days, the whole of the carbonate of lime was nearly dissolved. A series of membranes now only remained, of which the epidermis constituted the first. These membranes still retained the figure of the shell. The carbonate of lime was at first readily dissolved, because the acid came easily in contact with it; but the process became slower, as it was more difficult for the acid to insinuate itself between the different membranes of which the shell is composed. The haliotis iris, and the turbo olearius, were found to resemble this mussel, except that the membranaceous parts were more compact and dense.

When these shells are deprived by an acid of the carbonate of lime, which gives them their hardness, they appear to be formed of different membranes, applied stratum super stratum. Each membrane is furnished with a corresponding coat or crust of carbonate of lime, and it is so situated, that it is always between every two membranes, beginning with the epidermis, and ending with the internal membrane, which has been last formed. The animals which inhabit these stratified shells, increase their habitation by the addition of a stratum of carbonate of lime, which is secured by a new membrane. And as every additional stratum exceeds in extent that which was previously formed, the shell becomes stronger in proportion as it is enlarged; and thus the growth and age of the animal may be denoted by the number of strata of which the shell is composed. Similar experiments were made

on pieces of mother-of-pearl as they are imported from China, and with precisely the fame refults. They appeared to be composed of the fame gelatinous matter and carbonate of lime. In all the shells of this description which were immersed in acids, the membranaceous parts retained the exact figure of the shell, and they appeared distinctly to be composed of fibres, arranged in a parallel direction, corresponding to the configuration of the shell.

Pearl.—The constituent parts of pearl appear to be familiar to those of mother-of-pearl. They are composed of concentric coats of membrane and carbonate of lime, and resemble in structure the globular, calcareous concretions which are known by the name of pisolites. The iridescence and undulated appearance of pearl and mother-of-pearl, evidently depend on their lamellated structure and semitransparency.

From their experiments, it appears, that shells are composed of carbonate of lime and gluten. In those as pure in the porcellaneous shells, the proportion of carbonate of lime is great, while that of the animal matter is small; and these may be regarded as the beginning of the series; while shells that come under the description of mother-of-pearl are to be placed at the other extremity, having a smaller proportion of carbonate of lime, and a greater proportion of membranaceous fulblance. In the first the carbonate of lime is merely cemented by the animal matter; in the latter the carbonate of lime serves to harden the membranaceous fulblance. But between these two extremes, in the proportion of carbonate of lime and animal gluten, of which all tealaceous fulblances are composed, there are no doubt numerous intermediate gradations, arising from the nature of the animal to which they form a covering, its peculiar habits, or mode of life.

Sect. II. Of the Formation of Shells.

The shell or covering of testaceous animals, has been considered as in some measure analogous to the bones of other animals, although its formation and growth are very different, since it serves as a base or support to the muscles, which are attached to its internal surface. The principal use of the shell, however, is to serve as a covering or defence to the animal.

Telacous animals are not only extremely different in external form, but also in the mode of their production. Some are viviparous, as the most of those which inhabit bivalve shells, multivalves, and even some of the univalves; while the others, which form the far greater proportion, are oviparous. In one point, however, they all agree, that whatever be the mode of production, whether from an egg, or directly from the uterus of the mother, the shell is formed on the body of the young animal, and is proportioned to its bulk.

the young animal, and is proportioned to its
The best observations which have yet been
the most elaborate investigation which has

the most elaborate investigation which has hitherto ap-
peared, concerning the formation and development of
shells, are those of the celebrated Reaumur, which were
published in the Memoirs of the Academy of Sciences
for the year 1709. The fame subject has been prosecuted
by other authors, but their results have been nearly
the fame as those of this distinguished naturalist. Klein
is almost the only author who has advanced a different
opinion. In his dissertation concerning the formation


\section*{Source record}
\begin{description}
\item[Scan ID] \texttt{britannica:309f765b1ecc83e0a7f7762e44b4a927a7555443:train:473}
\item[Image file] \texttt{page-0473-4ab88a7103cd.jpg}
\item[Image SHA-256] \texttt{4ab88a7103cd8f6c0c4f49cde0370890e3cbca2af72a6a90f8e7db56d053c32d}
\item[Dataset] NationalLibraryOfScotland/encyclopaedia\_britannica\_illustrated at revision \texttt{309f765b1ecc83e0a7f7762e44b4a927a7555443} (CC0-1.0)
\item[Extraction] model \texttt{nvidia/nemotron-3.5-parse}, raw reply \texttt{extract/raw/page-0473-4ab88a7103cd.json}
\item[Raw reply SHA-256] \texttt{aaf9311cf965cbb3e6f28e5eaabfc9b9d30ea77f74f0b6bbe048d0b27022e228}
\item[Correction] model \texttt{nvidia/Nemotron-3.5-Super-VL}, block crops, raw replies \texttt{correct/raw/crop\_nothink/page-0473-4ab88a7103cd-b*.json}; super\_vl: 15
\item[Running header] none
\end{description}
\end{document}
